\chapter{Intrauterine Growth Restriction}
\index{Intrauterine Growth Restriction}
\label{sec:Intrauterine Growth Restriction}

\section{Overview}
Intrauterine growth restriction (IUGR), also termed fetal growth restriction (FGR), is defined as an estimated fetal weight below the 10th percentile for gestational age, often subclassified as symmetric (early onset, involving all biometric parameters) and asymmetric (late onset, with head sparing). It affects 5--10\% of pregnancies.
\section{Causes}
Causes include maternal factors (high blood pressure, preeclampsia, chronic kidney disease, malnutrition, smoking, alcohol, illicit drugs), placental factors (abnormal placentation, tissue death due to lack of blood supply, velamentous cord insertion), and fetal factors (aneuploidy, congenital infections, multiple gestation).     
\section{Risk Factors}

\begin{itemize}[noitemsep]
  \item Genetic predisposition and family history
  \item Advancing age
  \item Lifestyle factors (diet, exercise, smoking, alcohol)
  \item Environmental exposures
  \item Underlying medical conditions
  \item Medication use
\end{itemize}
\section{Signs and Symptoms}

\subsection{Common symptoms}
\begin{itemize}[noitemsep]
  \item \item You may experience subnormal fetal growth on serial ultrasound. Pain or discomfort in the affected area    \item Pain or discomfort in the affected area
  \end{itemize}

\subsection{Emergency warning signs}
\begin{itemize}[noitemsep]
  \item Severe or worsening symptoms of intrauterine growth restriction require immediate medical evaluation
\end{itemize}
\section{Progression and Stages}
The condition may progress through different stages of severity over time. Early stages often involve mild or subtle symptoms that may be overlooked. Moderate stages involve more noticeable symptoms affecting daily function. Advanced stages may involve significant impairment and complications.
\section{Complications}
\begin{itemize}[noitemsep]
  \item Disease progression leading to worsening symptoms
  \item Secondary infections or injuries
  \item Side effects of treatment
  \item Reduced quality of life
  \item Emotional and psychological impact
\end{itemize}
\section{Diagnosis}
\begin{itemize}[noitemsep]
  \item Diagnosis relies on accurate gestational dating and serial ultrasound measurement of fetal biometry, estimated fetal weight, and amniotic fluid volume
  \item Umbilical artery Doppler velocimetry identifies fetuses at highest risk.
  \item Comprehensive medical history and physical examination
  \item Laboratory tests to assess organ function
  \item Imaging studies to visualize affected structures
  \item Specialized testing depending on the affected system
  \item Consultation with relevant specialists
\end{itemize}
\section{Prevention}
\begin{itemize}[noitemsep]
  \item Maintain a healthy lifestyle with balanced diet and exercise
  \item Avoid known risk factors
  \item Attend regular medical check-ups for early detection
  \item Follow recommended screening guidelines
\end{itemize}
\section{Treatment Options}
Treatment options include antenatal surveillance and delivery when fetal maturity is achieved or when the risks of continued in utero existence outweigh neonatal risks. Medications to manage symptoms and address underlying mechanisms Lifestyle modifications and self-care measures Physical therapy and rehabilitation Surgical intervention when indicated Regular monitoring and follow-up care Patient education and support services

\begin{itemize}[noitemsep]
  \item Medications to manage symptoms and address underlying mechanisms
  \item Lifestyle modifications and self-care measures
  \item Physical therapy and rehabilitation
  \item Surgical intervention when indicated
  \item Regular monitoring and follow-up care
\end{itemize}
\section{Living with It}
\begin{itemize}[noitemsep]
  \item Follow your treatment plan as prescribed
  \item Maintain regular follow-up with your healthcare provider
  \item Make necessary lifestyle adjustments
  \item Seek support from family, friends, or support groups
  \item Stay informed about your condition
\end{itemize}
\section{Outlook}
Your outlook depends on the underlying cause and gestational age at delivery. The outlook varies depending on the specific condition, severity at diagnosis, and response to treatment. Early intervention generally leads to better outcomes. Many conditions can be effectively managed with appropriate treatment. Regular follow-up care is important for monitoring and treatment adjustment.
